\ifdefined\gkver\else\def\gkver{student}\fi
\documentclass[\gkver]{gaokaozhenti}
\begin{document}

\chapter{2001年上海理综}
%% paperType: 综合能力测试物理部分

\begin{enumerate}
\item
%% number: 13
%% paperName: 2001年上海理综第 13 题 3 分
%% typeId: 1
%% type: 单选题
%% chapter: 电场
%% point: 电场线
%% method: 无
%% score: 3
%% degree: 750
%% duplicateId: 0
%% body:
法拉第首先提出用电场线形象生动地描绘电场。图为点电荷 $a$、$b$ 所形成电场的电场线分布图，以下几种说法正确的是\xzanswer{B}

\onepicture[w=6.0cm]{figs/13.jpg}

\fourchoices[answer=B]
{$a$、$b$ 为异种电荷，$a$ 带电量大于 $b$ 带电量}
{$a$、$b$ 为异种电荷，$a$ 带电量小于 $b$ 带电量}
{$a$、$b$ 为同种电荷，$a$ 带电量大于 $b$ 带电量}
{$a$、$b$ 为同种电荷，$a$ 带电量小于 $b$ 带电量}

%% answer: B
%% memo:
\memoanswer{
本题原卷及材料中未提供详解，待补充。
}

\item
%% number: 14
%% paperName: 2001年上海理综第 14 题 3 分
%% typeId: 1
%% type: 单选题
%% chapter: 电磁感应
%% point: 楞次定律
%% method: 无
%% score: 3
%% degree: 650
%% duplicateId: 0
%% body:
某实验小组用如图所示的实验装置来验证楞次定律。当条形磁铁自上而下穿过固定的线圈时，通过电流计的感应电流方向是\xzanswer{D}

\onepicture[w=5.0cm]{figs/14.png}

\fourchoices[answer=D]
{$a\to G\to b$}
{先 $a\to G\to b$，后 $b\to G\to a$}
{$b\to G\to a$}
{先 $b\to G\to a$，后 $a\to G\to b$}

%% answer: D
%% memo:
\memoanswer{
本题原卷及材料中未提供详解，待补充。
}

\item
%% number: 15
%% paperName: 2001年上海理综第 15 题 3 分
%% typeId: 1
%% type: 单选题
%% chapter: 电路
%% point: 串并联电路
%% method: 无
%% score: 3
%% degree: 550
%% duplicateId: 0
%% body:
某实验小组用三只相同的小灯泡，联成如图所示的电路，研究串并联电路特点。实验中观察到的现象是\xzanswer{D}

\onepicture[w=6.0cm]{figs/15.jpg}

\fourchoices[answer=D]
{$K_{2}$ 断开，$K_{1}$ 与 $a$ 连接，三只灯泡都熄灭}
{$K_{2}$ 断开，$K_{1}$ 与 $b$ 连接，三只灯泡亮度相同}
{$K_{2}$ 闭合，$K_{1}$ 与 $a$ 连接，三只灯泡都发光，$L_{1}$、$L_{2}$ 亮度相同}
{$K_{2}$ 闭合，$K_{1}$ 与 $b$ 连接，三只灯泡都发光，$L_{3}$ 亮度小于 $L_{2}$ 的亮度}

%% answer: D
%% memo:
\memoanswer{
本题原卷及材料中未提供详解，待补充。
}

\item
%% number: 16
%% paperName: 2001年上海理综第 16 题 3 分
%% typeId: 1
%% type: 单选题
%% chapter: 气体性质
%% point: 查理定律
%% method: 无
%% score: 3
%% degree: 750
%% duplicateId: 0
%% body:
在冬季，剩有半瓶热水的暖水瓶经过一个夜晚后，第二天披瓶口的软木塞时觉得很紧，不易拔出来。其中主要原因是\xzanswer{D}

\fourchoices[answer=D]
{软木塞受潮膨胀}
{瓶口因温度降低而收缩变小}
{白天气温升高，大气压强变大}
{瓶内气体因温度降低而压强减小}

%% answer: D
%% memo:
\memoanswer{
本题原卷及材料中未提供详解，待补充。
}

\item
%% number: 17
%% paperName: 2001年上海理综第 17 题 3 分
%% typeId: 1
%% type: 单选题
%% chapter: 运动和力
%% point: 变加速直线运动
%% method: 无
%% score: 3
%% degree: 550
%% duplicateId: 0
%% body:
在一种叫做“蹦极跳”的运动中，质量为 $m$ 的游戏者身系一根长为 $L$、弹性优良的轻质柔软橡皮绳，从高处由静止开始下落 $1.5L$ 时到达最低点，若在下落过程中不计空气阻力，则以下说法正确的是\xzanswer{A}

\fourchoices[answer=A]
{速度先增大后减小}
{加速度先减小后增大}
{动能增加了 $mgL$}
{重力势能减少了 $mgL$}

%% answer: A
%% memo:
\memoanswer{
本题原卷及材料中未提供详解，待补充。
}

\item
%% number: 26
%% paperName: 2001年上海理综第 26 题 6 分
%% typeId: 3
%% type: 填空题
%% chapter: 运动和力
%% point: 牛顿定律的应用
%% method: 无
%% score: 6
%% degree: 650
%% duplicateId: 0
%% body:
1999 年 11 月 20 日我国成功发射和回收了“神舟”号实验飞船，标志着我国的运载火箭技术水平已跻身于世界先进行列。

图中 $A$ 为某火箭发射场，$B$ 为山区，$C$ 为城市。发射场正在进行某型号火箭的发射试验。该火箭起飞时质量为 $2.02\times10^{5}\Ukg$，起飞推力 $2.75\times10^{6}\UN$，火箭发射塔高 $100\Um$，则该火箭起飞时的加速度大小为 \tkanswer[2cm]{$3.81$}$\Umsq$；在火箭推力不变的情况下，若不考虑空气阻力及火箭质量的变化，火箭起飞后，经 \tkanswer[2cm]{$7.25$}$\Us$ 飞离火箭发射塔。（参考公式及常数：$F_{\text{合}}=ma$，$v_{t}=v_{0}+at$，$s=v_{0}t+\frac{1}{2}at^{2}$，$g=9.8\Umsq$）

\onepicture[w=8.5cm, align=right]{figs/26.jpg}

%% answer: 
%% memo:
\memoanswer{
本题原卷及材料中未提供详解，待补充。

（说明：原 PDF 题干印作“起飞时质量为 $2.02\times10^{3}\Ukg$”，与参考答案 $3.81\Umsq$、$7.25\Us$ 不自洽。由 $a=\frac{F}{m}-g$、$t=\sqrt{\frac{2s}{a}}$ 反推，质量应为 $2.02\times10^{5}\Ukg$，故据此径改，原值备考。）
}

\item
%% number: 27
%% paperName: 2001年上海理综第 27 题 4 分
%% typeId: 3
%% type: 填空题
%% chapter: 光学
%% point: 光的电磁说
%% method: 无
%% score: 4
%% degree: 650
%% duplicateId: 0
%% body:
为了转播火箭发射的实况，在发射场建立了发射台用于发射广播与电视信号。已知传输无线电广播所用的电磁波波长为 $550\Um$，而传输电视信号所用的电磁波波长为 $0.566\Um$，为了不让山区挡住信号的传播，使城市居民能收听和收看火箭发射的实况，必须通过建在山顶上的转发站来转发 \tkanswer[1.8cm]{电视信号}（选填：无线电广播信号或电视信号）。这是因为：\tkanswer[8cm]{电视信号的波长短，沿直线传播，受山坡阻挡，不易衍射。}

%% answer: 
%% memo:
\memoanswer{
本题原卷及材料中未提供详解，待补充。
}

\item
%% number: 33
%% paperName: 2001年上海理综第 33 题 4 分
%% typeId: 3
%% type: 填空题
%% chapter: 功和能
%% point: 机械能守恒定律
%% method: 无
%% score: 4
%% degree: 650
%% duplicateId: 0
%% body:
随着人类能量消耗的迅速增加，如何有效地提高能量利用率是人类所面临的一项重要任务。下图是上海“明珠线”某轻轨车站的设计方案，与站台连接的轨道有一个小的坡度。请你从提高能量利用效率的角度，分析这种设计的优点。

\onepicture[w=13.0cm, align=right]{figs/33.jpg}

%% answer: 
\jdanswer{列车进站时，利用上坡使部分动能转化为重力势能，减少因为刹车而损耗的机械能；列车出站时利用下坡把储存的重力势能又转化为动能，起到节能作用。（若从力的角度分析，得 2 分）}

%% memo:
\memoanswer{
本题原卷及材料中未提供详解，待补充。
}

\item
%% number: 39
%% paperName: 2001年上海理综第 39 题 3 分
%% typeId: 3
%% type: 填空题
%% chapter: 电路
%% point: 电功电功率
%% method: 无
%% score: 3
%% degree: 650
%% duplicateId: 0
%% body:
下图是两个电池外壳上的说明文字。

\onepicture[w=1.4cm]{figs/39.jpg}

\begin{center}
\footnotesize
\begin{tblr}{|c|c|}
\hline
某型号进口电池 & 某型号国产电池 \\
\hline
RECHARGEABLE & GNY 0.6（KR-AA） \\
\hline
$1.2\UV$，$500\UmAh$ & $1.2\UV$，$600\UmAh$ \\
\hline
STANDARD & RECHARGEBLE \\
\hline
CHARGE & STANDARD CHARGE \\
\hline
$15\Uh$ at $50\UmA$ & $15\Uh$ at $60\UmA$ \\
\hline
\end{tblr}
\end{center}

上述进口电池的电动势是 \tkanswer[1.5cm]{$1.2$}$\UV$。上述国产电池最多可放出 \tkanswer[2cm]{$600$}$\UmAh$ 的电量；若该电池平均工作电流为 $0.03\UA$，则最多可使用 \tkanswer[1.5cm]{$20$} 小时。

%% answer: 
%% memo:
\memoanswer{
本题原卷及材料中未提供详解，待补充。

（说明：原 PDF 表格中国产电池印作“$1.2\UV$，$60\UmAh$”，与参考答案 $600\UmAh$、$20$ 小时不自洽；由 $t=\frac{Q}{I}=\frac{600\UmAh}{30\UmA}=20\Uh$ 可知应为 $600\UmAh$，故据此径改，原值备考。）
}

\item
%% number: 46
%% paperName: 2001年上海理综第 46 题 1 分
%% typeId: 1
%% type: 单选题
%% chapter: 其他
%% point: 物理思想
%% method: 无
%% score: 1
%% degree: 850
%% duplicateId: 0
%% body:
人类在探索自然规律的进程中总结了许多科学方法。如分析归纳法、演绎法、等效替代法、控制变量法、理想实验法等。在下列研究中，运用理想实验方法进行研究的是\xzanswer{D}

\fourchoices[answer=D]
{爱因斯坦提出光子假说}
{查理得出气体状态变化的规律}
{卢瑟福提出原子的核式结构模型}
{伽利略得出力不是维持物体运动原因的结论}

%% answer: D
%% memo:
\memoanswer{
本题原卷及材料中未提供详解，待补充。
}

\end{enumerate}

\end{document}
